%% kapa: exam=1 round=24 year=2013 period=1 form=A subject=economics subject_ko=경제원론 seq=3 number=43
%% source: Q-Net 감정평가사 기출문제 (cst00309) artlSeq=5073802 file=extracted/1008/5073802/2013년 제24회 감정평가사 1차시험 기출문제/1교시(A형).hwp
%% conversion: kordoc 4.24.51 -> tools/gen1.py (auto); review=proofread-sonnet; flags=none
\begin{question}{43}
( )에 들어갈 내용을 순서대로 옳게 연결한 것은?
\begin{framedtext}
위험애호적(risk-loving)인 사람의 폰 노이만-모겐스턴 효용함수(von Neumann Morgenstern utility function)는 ( ㄱ )함수이며, ( ㄴ )이(가) 기대소득보다 크므로 위험 프리미엄이 0보다 ( ㄷ ).
\end{framedtext}
\begin{choices}
  \item ㄱ: 오목, ㄴ: 기대효용, ㄷ: 크다
  \item ㄱ: 볼록, ㄴ: 기대효용, ㄷ: 작다
  \item ㄱ: 오목, ㄴ: 기대효용, ㄷ: 작다
  \item ㄱ: 오목, ㄴ: 확실성등가, ㄷ: 크다
  \item ㄱ: 볼록, ㄴ: 확실성등가, ㄷ: 작다
\end{choices}

eviewflag{stem and choices 1-4 overlap/missing in rhwp render; only box and choice 5 verified}
\end{question}
\answer{5}
